% Options for packages loaded elsewhere
\PassOptionsToPackage{unicode}{hyperref}
\PassOptionsToPackage{hyphens}{url}
%
\documentclass[
]{article}
\usepackage{amsmath,amssymb}
\usepackage{lmodern}
\usepackage{iftex}
\ifPDFTeX
  \usepackage[T1]{fontenc}
  \usepackage[utf8]{inputenc}
  \usepackage{textcomp} % provide euro and other symbols
\else % if luatex or xetex
  \usepackage{unicode-math}
  \defaultfontfeatures{Scale=MatchLowercase}
  \defaultfontfeatures[\rmfamily]{Ligatures=TeX,Scale=1}
\fi
% Use upquote if available, for straight quotes in verbatim environments
\IfFileExists{upquote.sty}{\usepackage{upquote}}{}
\IfFileExists{microtype.sty}{% use microtype if available
  \usepackage[]{microtype}
  \UseMicrotypeSet[protrusion]{basicmath} % disable protrusion for tt fonts
}{}
\makeatletter
\@ifundefined{KOMAClassName}{% if non-KOMA class
  \IfFileExists{parskip.sty}{%
    \usepackage{parskip}
  }{% else
    \setlength{\parindent}{0pt}
    \setlength{\parskip}{6pt plus 2pt minus 1pt}}
}{% if KOMA class
  \KOMAoptions{parskip=half}}
\makeatother
\usepackage{xcolor}
\IfFileExists{xurl.sty}{\usepackage{xurl}}{} % add URL line breaks if available
\IfFileExists{bookmark.sty}{\usepackage{bookmark}}{\usepackage{hyperref}}
\hypersetup{
  pdftitle={Measuring Interdisciplinarity with Internal University Data: An Operational Guide},
  pdfauthor={A. Rivero},
  hidelinks,
  pdfcreator={LaTeX via pandoc}}
\urlstyle{same} % disable monospaced font for URLs
\usepackage{longtable,booktabs,array}
\usepackage{calc} % for calculating minipage widths
% Correct order of tables after \paragraph or \subparagraph
\usepackage{etoolbox}
\makeatletter
\patchcmd\longtable{\par}{\if@noskipsec\mbox{}\fi\par}{}{}
\makeatother
% Allow footnotes in longtable head/foot
\IfFileExists{footnotehyper.sty}{\usepackage{footnotehyper}}{\usepackage{footnote}}
\makesavenoteenv{longtable}
\setlength{\emergencystretch}{3em} % prevent overfull lines
\providecommand{\tightlist}{%
  \setlength{\itemsep}{0pt}\setlength{\parskip}{0pt}}
\setcounter{secnumdepth}{-\maxdimen} % remove section numbering
\ifLuaTeX
  \usepackage{selnolig}  % disable illegal ligatures
\fi

\title{Measuring Interdisciplinarity with Internal University Data: An
Operational Guide}
\author{A. Rivero}
\date{2026}

\begin{document}
\maketitle
\begin{abstract}
Universities possess rich internal data --- publication repositories,
co-authorship records, PhD supervision logs, and grant databases ---
that external bibliometric services cannot access. We present an
operational measurement protocol that leverages this institutional data
to assess interdisciplinary research at the departmental and
institutional level. The protocol combines a three-component
bibliometric panel (diversity, coherence, cross-field effect) with
institutional-only indicators (co-authorship diversity, co-supervision
diversity, grant panel diversity). Using a mock departmental scenario,
we illustrate how institutional data reveals polymathic breadth versus
genuine integration --- a distinction invisible to external databases.
We identify the main data obstacle (citation data for cross-field
effect) and propose a two-component workaround for institutions without
external database subscriptions.
\end{abstract}

\hypertarget{introduction}{%
\section{Introduction}\label{introduction}}

Universities increasingly seek to foster and assess interdisciplinary
research as part of strategic planning and quality assurance. However,
commercial bibliometric databases (Web of Science, Scopus) provide only
partial views of research activity, and their subscription costs can be
prohibitive. Moreover, these databases lack access to institutional
processes --- co-authorships, PhD co-supervisions, grant collaborations,
and departmental affiliations --- that reveal how knowledge integration
actually occurs within the university.

This note presents an operational guide for universities to measure
interdisciplinarity using primarily internal data. The core approach
combines a three-component bibliometric panel --- diversity
(\(\Delta\)), coherence (\(S\)), and cross-field effect (\(E\)),
developed in a companion review (Rivero, 2026) --- with
institutional-only indicators that capture collaborative structures. We
demonstrate the protocol on a mock departmental scenario and provide
implementation guidance for data extraction, quality checks, and
reporting.

\hypertarget{data-landscape}{%
\section{Data Landscape}\label{data-landscape}}

\hypertarget{what-universities-have}{%
\subsection{What Universities Have}\label{what-universities-have}}

Most research-intensive universities maintain:

\begin{itemize}
\tightlist
\item
  \textbf{Publication repositories}: Institutional repositories or CRIS
  (Current Research Information Systems) recording all faculty
  publications with metadata (authors, references, keywords).
\item
  \textbf{HR systems}: Departmental affiliations of all staff.
\item
  \textbf{PhD records}: Supervision records, including co-supervisors
  and their departmental affiliations.
\item
  \textbf{Grant databases}: Internal records of all grants, including
  funding source, review panel assignment, and co-investigators.
\end{itemize}

These data sources are comprehensive within the institution but
unavailable to external services.

\hypertarget{what-universities-need-but-may-lack}{%
\subsection{What Universities Need (But May
Lack)}\label{what-universities-need-but-may-lack}}

\begin{itemize}
\tightlist
\item
  \textbf{Disciplinary classification of references}: Mapping cited
  works to subject categories (e.g., Web of Science categories). Some
  repositories include this; most do not.
\item
  \textbf{Similarity matrix between categories}: Required for computing
  Rao-Stirling diversity. Can be derived from citation data or, more
  recently, estimated via large language models (Cantone, 2025).
\item
  \textbf{Citation data}: Who cites each publication, and from which
  fields. Essential for computing the cross-field effect (\(E\)).
  Available from Web of Science, Scopus, or the open-access OpenAlex,
  but not from internal systems.
\end{itemize}

The main data gap is citation data for computing \(E\). The protocol
below addresses this with a two-component workaround.

\hypertarget{measurement-protocol}{%
\section{Measurement Protocol}\label{measurement-protocol}}

\hypertarget{step-1-bibliometric-panel}{%
\subsection{Step 1: Bibliometric
Panel}\label{step-1-bibliometric-panel}}

For each researcher, compute:

\textbf{Diversity} (\(\Delta\)): The Rao-Stirling index over cited
references:
\[\Delta = \sum_{i \neq j} d_{ij}\, p_i\, p_j, \qquad d_{ij} = 1 - s_{ij}\]
where \(p_i\) is the proportion of references in category \(i\),
\(s_{ij}\) is the pairwise similarity between categories, and \(d_{ij}\)
is the corresponding distance. This requires reference classification
and a similarity matrix.

\textbf{Coherence} (\(S\)): The mean pairwise bibliographic coupling
among publications:
\[S = \frac{1}{\binom{n}{2}} \sum_{k < l} \cos(\mathbf{r}_k, \mathbf{r}_l)\]
where \(\mathbf{r}_k\) is the reference vector of publication \(k\).
This requires only internal repository data.

\textbf{Cross-field effect} (\(E\)): The fraction of citations received
from outside the researcher's primary category. This requires external
citation data.

\hypertarget{step-2-institutional-only-indicators}{%
\subsection{Step 2: Institutional-Only
Indicators}\label{step-2-institutional-only-indicators}}

Beyond the bibliometric panel, compute:

\textbf{Co-authorship diversity}: The fraction of publications with at
least one co-author from a different department:
\[\text{CoAuth} = \frac{\text{publications with cross-departmental co-authors}}{\text{total publications}}\]

\textbf{Co-supervision diversity}: The fraction of PhD students
co-supervised with faculty from other departments:
\[\text{CoSup} = \frac{\text{cross-departmental co-supervisions}}{\text{total supervisions}}\]

\textbf{Grant panel diversity}: The number of distinct funding review
panels from which the researcher has secured grants. This proxies
disciplinary breadth of funding capacity.

\hypertarget{step-3-interpretation}{%
\subsection{Step 3: Interpretation}\label{step-3-interpretation}}

The combination of bibliometric and institutional indicators
discriminates integration from polymathy:

\begin{longtable}[]{@{}
  >{\raggedright\arraybackslash}p{(\columnwidth - 10\tabcolsep) * \real{0.1667}}
  >{\raggedright\arraybackslash}p{(\columnwidth - 10\tabcolsep) * \real{0.1296}}
  >{\raggedright\arraybackslash}p{(\columnwidth - 10\tabcolsep) * \real{0.1296}}
  >{\raggedright\arraybackslash}p{(\columnwidth - 10\tabcolsep) * \real{0.1481}}
  >{\raggedright\arraybackslash}p{(\columnwidth - 10\tabcolsep) * \real{0.1296}}
  >{\raggedright\arraybackslash}p{(\columnwidth - 10\tabcolsep) * \real{0.2963}}@{}}
\toprule
\begin{minipage}[b]{\linewidth}\raggedright
Profile
\end{minipage} & \begin{minipage}[b]{\linewidth}\raggedright
\(\Delta\)
\end{minipage} & \begin{minipage}[b]{\linewidth}\raggedright
S
\end{minipage} & \begin{minipage}[b]{\linewidth}\raggedright
CoAuth
\end{minipage} & \begin{minipage}[b]{\linewidth}\raggedright
CoSup
\end{minipage} & \begin{minipage}[b]{\linewidth}\raggedright
Interpretation
\end{minipage} \\
\midrule
\endhead
Integrator & High & Moderate-high & High & High & Cross-disciplinary in
publications AND processes \\
Polymath & High & Low & Low & Low & Broad references but no
collaborative integration \\
Specialist & Low & High & Low & Low & Focused, disciplinary
researcher \\
\bottomrule
\end{longtable}

\textbf{Key insight:} A researcher with high \(\Delta\) and high grant
diversity might appear interdisciplinary from external data, but zero
co-authorship and co-supervision diversity reveal polymathic breadth
without integration. Institutional data provides this discriminatory
power.

This distinction mirrors the input/output separation emphasized in the
companion review: high diversity of inputs (\(\Delta\)) can indicate
either multidisciplinary breadth or integrative interdisciplinarity, and
only coherence/process indicators can separate the two reliably.

\hypertarget{demonstration-mock-department}{%
\subsection{Demonstration: Mock
Department}\label{demonstration-mock-department}}

A small Physics \& Materials Science department with 3 researchers
illustrates the protocol. We use five Web of Science-style categories
and an illustrative similarity matrix:

\begin{longtable}[]{@{}ll@{}}
\toprule
ID & Category \\
\midrule
\endhead
C1 & Physics, condensed matter \\
C2 & Materials science \\
C3 & Chemistry, physical \\
C4 & Optics \\
C5 & Engineering, electrical \\
\bottomrule
\end{longtable}

\begin{longtable}[]{@{}lrrrrr@{}}
\toprule
& C1 & C2 & C3 & C4 & C5 \\
\midrule
\endhead
C1 & 1.00 & 0.60 & 0.40 & 0.35 & 0.30 \\
C2 & 0.60 & 1.00 & 0.50 & 0.25 & 0.40 \\
C3 & 0.40 & 0.50 & 1.00 & 0.30 & 0.20 \\
C4 & 0.35 & 0.25 & 0.30 & 1.00 & 0.45 \\
C5 & 0.30 & 0.40 & 0.20 & 0.45 & 1.00 \\
\bottomrule
\end{longtable}

Similarity values are illustrative; in practice they would be derived
from inter-category citation patterns or estimated via large language
models (Cantone, 2025).

All coherence values (\(S\)) below are illustrative; per-publication
reference vectors are omitted for brevity.

\textbf{Dr.~Emma} (10 years post-PhD,
\(\mathbf{p}_E = (0.40, 0.30, 0.25, 0.03, 0.02)\)): - Biblio panel:
\(\Delta = 0.42\), \(S = 0.55\), \(E = 0.22\) - Institutional: CoAuth =
0.20, CoSup = 0.20, Grants = 2 panels - \textbf{Profile}: Integrator.
Moderate bibliometric diversity reinforced by cross-departmental
collaborations and co-supervisions.

\textbf{Dr.~Farid} (7 years post-PhD,
\(\mathbf{p}_F = (0.25, 0.25, 0.25, 0.20, 0.05)\)): - Biblio panel:
\(\Delta = 0.58\), \(S = 0.05\), \(E = 0.08\) - Institutional: CoAuth =
0.00, CoSup = 0.00, Grants = 4 panels - \textbf{Profile}: Polymath. High
diversity and grant breadth, but zero collaborative integration. Each
paper is a disconnected single-field contribution.

\textbf{Dr.~Greta} (4 years post-PhD,
\(\mathbf{p}_G = (0.70, 0.25, 0.03, 0.01, 0.01)\)): - Biblio panel:
\(\Delta = 0.28\), \(S = 0.75\), \(E = 0.08\) - Institutional: CoAuth =
0.00, CoSup = N/A, Grants = 1 panel - \textbf{Profile}: Early-career
specialist. Focused research program; low \(E\) expected at this stage.

Dr.~Farid's case is instructive: his grant diversity (4 panels) might
suggest strong interdisciplinary capacity, but the institutional
indicators reveal this is breadth without depth. The university sees
what external databases cannot.

\hypertarget{implementation-guidance}{%
\section{Implementation Guidance}\label{implementation-guidance}}

\hypertarget{data-extraction}{%
\subsection{Data Extraction}\label{data-extraction}}

\begin{enumerate}
\def\labelenumi{\arabic{enumi}.}
\item
  \textbf{Publication data}: Export from institutional repository with
  full reference lists. If reference categories are not already
  assigned, they must be mapped to a disciplinary taxonomy (e.g., Web of
  Science categories). This is the main data cleaning task.
\item
  \textbf{Co-authorship affiliations}: Extract from HR system. Flag all
  publications where co-authors belong to different departments.
\item
  \textbf{PhD supervision records}: Extract from graduate office.
  Identify co-supervisions and co-supervisors' departmental
  affiliations.
\item
  \textbf{Grant records}: Extract from research office. Map each grant
  to the funding agency's review panel taxonomy.
\item
  \textbf{Citation data (if available)}: Import from OpenAlex (free),
  Web of Science, or Scopus. Match citing papers to categories for \(E\)
  computation.
\end{enumerate}

\hypertarget{quality-checks}{%
\subsection{Quality Checks}\label{quality-checks}}

\begin{itemize}
\tightlist
\item
  \textbf{Missing reference categories}: If \textgreater25\% of
  references lack category assignments, \(\Delta\) estimates will be
  unreliable (Nakhoda et al., 2023). Report coverage statistics.
\item
  \textbf{Early-career N/A values}: Researchers with \textless5 years
  post-PhD or \textless10 publications should have confidence intervals
  reported alongside point estimates.
\item
  \textbf{Threshold calibration and ambiguity}: The profile cutoffs used
  in this note are illustrative. Institutions should calibrate them to
  local distributions and mark cases whose confidence intervals cross
  category boundaries as ``ambiguous, needs qualitative review.''
\item
  \textbf{Departmental affiliation changes}: Researchers who have
  changed departments mid-career may have artificially inflated
  co-authorship diversity. Stratify by time period if needed.
\end{itemize}

\hypertarget{reporting}{%
\subsection{Reporting}\label{reporting}}

Present indicators at three levels:

\begin{enumerate}
\def\labelenumi{\arabic{enumi}.}
\tightlist
\item
  \textbf{Individual}: Provide each researcher with their profile
  (\(\Delta\), S, E, CoAuth, CoSup, Grants). This supports
  self-assessment and career planning.
\item
  \textbf{Departmental}: Report distributions and medians across the
  department. Identify outliers for strategic discussion.
\item
  \textbf{Institutional}: Aggregate across departments for
  university-wide strategic planning. Compare distributions across
  fields (e.g., STEM vs. Humanities).
\end{enumerate}

\textbf{Do not rank researchers by a single composite score.} The
indicators are multidimensional by design. Ranking collapses the
information and incentivizes gaming (Rafols, 2019).

\hypertarget{workaround-when-citation-data-unavailable}{%
\subsection{Workaround When Citation Data
Unavailable}\label{workaround-when-citation-data-unavailable}}

If the institution lacks access to citation databases, compute a
\textbf{two-component internal panel} (\(\Delta\), S) plus institutional
indicators (CoAuth, CoSup, Grants). This still discriminates integrators
from polymaths:

\begin{longtable}[]{@{}
  >{\raggedright\arraybackslash}p{(\columnwidth - 10\tabcolsep) * \real{0.1667}}
  >{\raggedright\arraybackslash}p{(\columnwidth - 10\tabcolsep) * \real{0.1296}}
  >{\raggedright\arraybackslash}p{(\columnwidth - 10\tabcolsep) * \real{0.1296}}
  >{\raggedright\arraybackslash}p{(\columnwidth - 10\tabcolsep) * \real{0.1481}}
  >{\raggedright\arraybackslash}p{(\columnwidth - 10\tabcolsep) * \real{0.1296}}
  >{\raggedright\arraybackslash}p{(\columnwidth - 10\tabcolsep) * \real{0.2963}}@{}}
\toprule
\begin{minipage}[b]{\linewidth}\raggedright
Profile
\end{minipage} & \begin{minipage}[b]{\linewidth}\raggedright
\(\Delta\)
\end{minipage} & \begin{minipage}[b]{\linewidth}\raggedright
S
\end{minipage} & \begin{minipage}[b]{\linewidth}\raggedright
CoAuth
\end{minipage} & \begin{minipage}[b]{\linewidth}\raggedright
CoSup
\end{minipage} & \begin{minipage}[b]{\linewidth}\raggedright
Interpretation
\end{minipage} \\
\midrule
\endhead
Integrator & High & Moderate-high & High & High & Integration via
references AND processes \\
Polymath & High & Low & Low & Low & Broad but disconnected \\
Specialist & Low & High & Low & Low & Focused, disciplinary \\
\bottomrule
\end{longtable}

The loss of \(E\) reduces discrimination power for assessing cross-field
\emph{impact}, but integration \emph{capacity} is still captured.

\hypertarget{limitations-and-extensions}{%
\section{Limitations and Extensions}\label{limitations-and-extensions}}

The protocol has several limitations. First, it assumes the institution
has clean, structured data. Many universities' CRIS systems are
incomplete or inconsistent. Data cleaning is a non-trivial prerequisite.
Second, the protocol does not address quality --- it characterizes the
type of interdisciplinarity, not whether it is good. Quality assessment
requires independent evaluation (peer review, citation percentiles,
etc.). Third, the classification thresholds (e.g., ``high''
\(\Delta \geq 0.40\)) are illustrative and should be calibrated against
the institution's empirical distributions before operational use.
Fourth, confidence intervals can be wide for individual-level profiles
when publication counts are small; operational classification should
therefore use interval-aware rules rather than point estimates alone.

Extensions could include text-based indicators (semantic similarity of
titles/abstracts across a researcher's portfolio), teaching-based
indicators (cross-departmental teaching assignments), and temporal
analysis (tracking how \(\Delta\), \(S\), \(E\) evolve over a
researcher's career).

\hypertarget{conclusions}{%
\section{Conclusions}\label{conclusions}}

Universities possess data that external bibliometric services cannot
see. Leveraging this internal data --- co-authorships, PhD
co-supervisions, grant collaborations --- provides discriminatory power
that purely bibliometric measures lack. The measurement protocol
presented here combines a three-component bibliometric panel with
institutional-only indicators to distinguish genuine cross-disciplinary
integration from polymathic breadth. The main implementation obstacle is
citation data for computing cross-field effect; a two-component
workaround (diversity + coherence + institutional indicators) provides
substantial discrimination power without external dependencies.
Institutional deployment requires clean data, careful quality checks,
and resistance to the temptation to reduce multidimensional profiles to
single-number rankings.

\hypertarget{references}{%
\section{References}\label{references}}

\begin{itemize}
\tightlist
\item
  Cantone, G. G. (2025). Estimation of disciplinary similarity with
  large language models. \emph{Scientometrics}, 130(10):5345--5373.
\item
  Nakhoda, M., Whigham, P., and Zwanenburg, S. (2023). Quantifying and
  addressing uncertainty in the measurement of interdisciplinarity.
  \emph{Scientometrics}, 128:6107--6127.
\item
  Rafols, I. (2019). S\&T indicators in the wild: contextualization and
  participation for responsible metrics. \emph{Research Evaluation},
  28(1):7--22.
\item
  Rivero, A. (2026). Measuring interdisciplinarity: A multi-component
  indicator panel for research evaluation. Manuscript.
\end{itemize}

\end{document}
